\section{Should Fabrication Start Now?}

\recbox{
\textbf{Yes.} The theoretical model has been through defect-finding,
margin-aware re-verification, hand-cross-checking, cross-validation against
a competing published approach, seed-robustness testing, and now two
independently-computed safety standards (worst-case and gaze-weighted)
that agree on which designs are safe. Further theoretical iteration has
reached diminishing returns; a physical prototype is the highest-value
next step.
}

\subsection{Which Design to Build First}

The gaze-weighted result changes the practical fabrication plan for the
better: \textbf{one family (hexagonal), one pitch (4000\,$\mu$m), covering
M=1 through M=8 with only the hole diameter changing between inserts}
(1250, 750, 600, 500, 400\,$\mu$m). This is simpler to fabricate than the
mixed-family worst-case-optimal set (\texttt{ring\_offset\_concentric} +
\texttt{hexagonal} + \texttt{square} across different M ranges), and every
design in it achieves genuine 0.00\% ghosting under realistic use, not
just a structural guarantee at one gaze angle.

\textbf{Recommended build order:}
\begin{enumerate}
  \item \textbf{Hexagonal, pitch=4000\,$\mu$m, N=20, a\_hole=400\,$\mu$m}
  (M=8 insert) --- smallest hole, test cutting-edge quality first, per the
  established fabrication discipline.
  \item \textbf{Hexagonal, pitch=4000\,$\mu$m, N=20, a\_hole=750\,$\mu$m}
  (M=2--3 insert) --- mid-range, a second data point on the same base
  frame.
  \item \textbf{Square, pitch=4000\,$\mu$m, N=20, a\_hole=2000\,$\mu$m}
  (M=0/global) --- confirms the model at the opposite end of the aperture
  range, different family for a cross-check.
\end{enumerate}

\subsection{What to Measure}
Per the project's established phone-camera afocal protocol: blur at each
design's own M setting, and ghosting specifically at gaze angles beyond
17$^\circ$ (the 95\%-of-time boundary) --- this is the one region where
the worst-case and gaze-weighted models could plausibly diverge from
reality in different ways, and it is the most informative test of which
safety standard the real optics actually respect.

\end{document}
